\documentclass[12pt,a4paper,landscape]{article}
\usepackage[spanish,es-nodecimaldot]{babel}
\usepackage[utf8]{inputenc}
\usepackage[T1]{fontenc}
\usepackage{lmodern}
\usepackage[top=1.1cm,bottom=1.1cm,left=1.4cm,right=1.4cm]{geometry}
\usepackage{amsmath,amssymb}
\usepackage[table]{xcolor}
\usepackage{booktabs,array}
\usepackage{enumitem}
\usepackage[normalem]{ulem}
\usepackage{multicol}

\setlength{\parindent}{0pt}
\setlength{\parskip}{7pt}
\setlength{\columnsep}{1.2cm}
\pagestyle{empty}

\newcommand{\cuerpo}{\fontsize{12.6}{15.5}\selectfont}
\newenvironment{sangria}{\begin{list}{}{\setlength{\leftmargin}{16pt}\setlength{\rightmargin}{0pt}\setlength{\itemindent}{0pt}\setlength{\listparindent}{0pt}\setlength{\topsep}{2pt}\setlength{\parsep}{4pt}}\item[]}{\end{list}}

\AtBeginDocument{%
  \setlength{\abovedisplayskip}{7pt}%
  \setlength{\belowdisplayskip}{7pt}%
  \setlength{\abovedisplayshortskip}{4pt}%
  \setlength{\belowdisplayshortskip}{4pt}%
  \cuerpo
}

\begin{document}

\begin{center}
{\fontsize{17}{21}\selectfont\bfseries\uline{Pruebas de aleatoreidad}} \quad --- Prueba de corridas
\end{center}

\vspace{2pt}

\begin{multicols}{2}

Sirven para comprobar si el orden de una muestra o proceso puede considerarse \textbf{aleatorio}. Se analiza si los valores se suceden sin un patrón sistemático.\\
Una \textbf{corrida} o \textbf{racha} es una secuencia de una o más letras iguales consecutivas.
\[
a\,a \mid b\,b \mid a \mid b\,b\,b \qquad \longrightarrow \quad \textbf{4 corridas}
\]

{\fontsize{15}{18}\selectfont\bfseries Hipotesis}

H0:el arreglo o secuencia es aleatorio

H1:el arreglo no es aleatorio $\Rightarrow$prueba bilateral\\
H1:hay demasiadas corridas $\Rightarrow$prueba unilateral derecha\\
H1:hay muy pocas corridas $\Rightarrow$prueba unilateral izquierda

{\fontsize{15}{18}\selectfont\bfseries Pasos}

\textbf{1. Construir la cadena de letras}

\begin{sangria}
-Si los datos son \textbf{binarios} ---éxito/fracaso, 1/0, caro/barato---, se construye directamente una cadena de letras a y b.\\
-Si los datos son \textbf{numéricos}:
\begin{enumerate}[leftmargin=20pt,itemsep=2pt,topsep=3pt,parsep=0pt]
  \item Calcular la mediana.
  \item Asignar a a cada observación superior a la mediana.
  \item Asignar b a cada observación inferior a la mediana.
  \item Los valores iguales a la mediana se omiten.
  \item Mantener el orden original de los datos.
\end{enumerate}
\end{sangria}

\begin{minipage}{\linewidth}
\textbf{2. Contar las cantidades}

\begin{sangria}
n1=cantidad de letras a\\
n2=cantidad de letras b\\
u=cantidad total de corridas
\end{sangria}
\end{minipage}

\textbf{3. Calcular los parámetros de la distribución de corridas}
\[
\mu_u = \frac{2n_1 n_2}{n_1+n_2} + 1
\]
\[
\sigma_u = \sqrt{\frac{2n_1 n_2\,(2n_1 n_2 - n_1 - n_2)}{(n_1+n_2)^2\,(n_1+n_2-1)}}
\]

\textbf{4. Aproximacion normal}

\begin{sangria}
si n1$>$10 y n2$>$10 la dist. del num de corridas \textbf{u} se aproxima a una dist. normal, aplico estadistico z
\[
Z_{\text{calc}} = \frac{u - \mu_u}{\sigma_u}
\]
\end{sangria}

\textbf{Interpretaciones}

\begin{itemize}[leftmargin=20pt,itemsep=4pt,topsep=3pt,parsep=0pt]
  \item {\small $u$ grande: demasiados cambios entre $a$ y $b$; alternancia excesiva.}
  \item {\small $u$ chico: pocas alternancias; los valores iguales tienden a agruparse.}
  \item {\small $u$ cercano a $\mu_u$: comportamiento compatible con aleatoriedad.}
\end{itemize}

\textbf{Se usa en control de calidad para dar aleatoreidad a procesos}

\vspace{4pt}
\begin{center}
{\small
\setlength{\tabcolsep}{4.2pt}
\begin{tabular}{l c c c c c c c c c c}
\toprule
\textbf{Dato} & .261 & .258 & .249 & .251 & .247 & .256 & .250 & .247 & .255 & .243 \\
\textbf{Signo} & a & a & b & a & b & a & {\footnotesize (Empate)} & b & a & b \\
\bottomrule
\end{tabular}
}
\end{center}

$H_0$ : la secuencia es aleatoria

$H_1$ : hay demasiadas corridas $(u > \mu_u)$

\end{multicols}

\end{document}
